\scanpage{137}\mspage{11}
% archive sheet number: 136
\deleted{\deleted{$\Phi$ une} partie \inserted{$\Phi$} de $\mathrm{Fl}\,M$ telle que
\[
C_0\subset\Phi\subset\widetilde C_0
\]
et telle que l'on ait ceci : $\Phi$ est \deleted{\unsure{stable par cochgt}} \inserted{\unsure{formé de flèches}} \deleted{\unsure{co-quarrables}}, et stable par cochangement de base, par sommes amalgamées infinies (cf. 1.18. p. 8), par limites ordinales strictes (cf. 1.19 p. 8). Si $\Phi$ est \unsure{choisie|choisi} ainsi on a ceci :}
\margin{Posons
\[
\left\{\begin{aligned}
\Phi&=\widetilde C_0=\Psi^*\\
\Psi&=C_{0*}=\Phi_*
\end{aligned}\right.\tag*{(1.23)}
\]}
Ceci posé, on a :

a) \deleted{\unsure{Pour tout}} \inserted{Pour} \unsure{tt} ordinal $\alpha$, on a $i_\alpha(f)\in\Phi\deleted{{}=\unsure{(C_{0*})^*|(C_0)^*}}$

b) Supposons que les sources des $i\in C_0$ soient accessibles, et soit donc $\unsure{\pi|I}$ un cardinal $\in\unsure{\mathcal U|U}$ tel qu'elles soient $\pi$-accessibles. Alors pour tout nombre ordinal $\alpha$ grand devant $\pi$, on a $p_\alpha(f)\in C_{0*}$, donc
\[
f=p_\alpha(f)i_\alpha(f)\quad\text{avec}\quad i_\alpha(f)\in\Phi,\quad p_\alpha(f)\in\Psi.
\]
On a un résultat légèrement plus précis, \deleted{que} que voici

\uline{Corollaire 1} \deleted{\unsure{Soit}} Soit $\Phi_0$ une partie de $\mathrm{Fl}\,M$ telle que
\[
C'''_0\subset\Phi_0\subset\Phi,\quad\text{(d'où}\quad C_0\subset\Phi_0\subset\Phi\text{)}\tag*{(1.24)}
\]
où $\unsure{\in|c}C'''_0$ \inserted{$\unsure{\subset|\in}\Phi$ est} \unsure{défini} dans 1.19. p. 8. Par exemple, si $\Phi_0$ est une partie de $\mathrm{Fl}\,M$ \deleted{\unsure{contenant}} satisfaisant \deleted{\unsure{telle}}
\[
C_0\subset\Phi_0\subset\Phi,
\]
et stable par \unsure{cochgt.} de base, par \inserted{petites} sommes amalgamées (cf. 1.18 p. 8)), par $\varinjlim$ ordinales strictes (cf. 1.19. p. 8), elle satisfait 1.25. (NB \unsure{Si} $M$ \unsure{admet} des sommes amalgamées simples, de tels $\Phi_0$ existent trivialement, il suffit de prendre $\Phi_0=\Phi$)
